\documentclass[../../../main.tex]{subfiles}

\begin{document}

\section{Isogonal Conjugate}

The isogonal conjugate of a point $P$ with respect to a triangle
is an interesting well-defined point that can be expanded to
different techniques and patterns. Although the name suggests
the definition of the point, here is the definition.

\begin{definition}{}{}
The \textbf{isogonal conjugate} of a point $P$ with respect to a
triangle $ABC$ is the intersection of the reflections of $AP, BP, CP$
with respect to the angle bisectors at $A$, $B$, and $C$.
\end{definition}

To better understand this topic, consider the following diagram.

\begin{figure}[H]
\centering
\begin{tikzpicture}[scale=1.2]
    \coordinate (A) at (4, 4);
    \coordinate (B) at (0, 0);
    \coordinate (C) at (6, 0);
    \coordinate (P) at (3,2);
    \coordinate (Q) at (4.444444444444444,0.8888888888888894);
    \coordinate (D) at (3,-2);
    \coordinate (E) at (6.2,3.6);
    \coordinate (F) at (2,3);

    \draw[thick] (A) -- (B) -- (C) -- (A);
    \draw[thick, BrickRed, dotted] (A) -- (P);
    \draw[thick, BrickRed, dotted] (B) -- (P);
    \draw[thick, BrickRed, dotted] (C) -- (P);
    \draw[thick, RoyalPurple, dotted] (A) -- (Q);
    \draw[thick, RoyalPurple, dotted] (B) -- (Q);
    \draw[thick, RoyalPurple, dotted] (C) -- (Q);

    \node[above] at (A) {\footnotesize $A$};
    \node[left] at (B) {\footnotesize $B$};
    \node[right] at (C) {\footnotesize $C$};
    \draw[fill=BrickRed] (P) circle (0.05)
        node[left] {\footnotesize $P$};
    \draw[fill=RoyalPurple] (Q) circle (0.05)
        node[below] {\footnotesize $Q$};

    \pic[draw,"$\bullet$", angle radius=8mm, angle eccentricity=1.4]
        {angle=B--A--P};
    \pic[draw,"$\bullet$", angle radius=8mm, angle eccentricity=1.4]
        {angle=Q--A--C};
    \pic[draw,"$\circ$", angle radius=8mm, angle eccentricity=1.4]
        {angle=P--B--A};
    \pic[draw,"$\circ$", angle radius=8mm, angle eccentricity=1.4]
        {angle=C--B--Q};
    \pic[draw,"$\times$", angle radius=8mm, angle eccentricity=1.4]
        {angle=Q--C--B};
    \pic[draw,"$\times$", angle radius=8mm, angle eccentricity=1.4]
        {angle=A--C--P};
\end{tikzpicture}
\end{figure}

Note that the lines $AQ, BQ, CQ$ are the \textbf{isogonal lines} of
$AP, BP, CP$ respectively. The point $Q$ is the isogonal conjugate
of $P$ with respect to $\triangle{ABC}$. Now this point $Q$ is
well-defined. To show that the isogonal conjugate is always a
well-defined point, we can use many different methods, but here
are two of my favorite ones.

For the first proof, we will use Angle-Ceva Theorem. I honestly am
not sure the standard term for the theorem because it is not
widely known, but in Korea we call it as Angle-Ceva Theorem and
I also saw some English forms like the angle form of Ceva' Theorem
and trigonometric form of Ceva's Theorem.

\begin{important}
The theorem states that
the following equation holds for some point $P$ in triangle $ABC$.
\[
\frac{\sin\angle{ABP}}{\sin\angle{PBC}} \cdot
    \frac{\sin\angle{CAP}}{\sin\angle{PAB}} \cdot
    \frac{\sin\angle{BCP}}{\sin\angle{PCA}}
= 1
\]
\end{important}

Note that by construction, the following equation holds.
\begin{align*}
    &\quad\ \frac{\sin\angle{CBQ}}{\sin\angle{QBA}} \cdot
        \frac{\sin\angle{ACQ}}{\sin\angle{QCB}} \cdot
        \frac{\sin\angle{BAQ}}{\sin\angle{QAC}} \\
    &= \frac{\sin\angle{ABP}}{\sin\angle{PBC}} \cdot
        \frac{\sin\angle{BCP}}{\sin\angle{PCA}} \cdot
        \frac{\sin\angle{CAP}}{\sin\angle{PAB}}
    = 1
\end{align*}
Therefore by reverse Angle-Ceva, $Q$ is well-defined.

Another way of proving the existence is by using reflections.

\begin{figure}[H]
\centering
\begin{tikzpicture}[scale=0.8]
    \coordinate (A) at (4, 4);
    \coordinate (B) at (0, 0);
    \coordinate (C) at (6, 0);
    \coordinate (P) at (3,2);
    \coordinate (Q) at (4.444444444444444,0.8888888888888894);
    \coordinate (D) at (3,-2);
    \coordinate (E) at (6.2,3.6);
    \coordinate (F) at (2,3);

    \draw[thick] (A) -- (B) -- (C) -- (A);
    \draw[thick, BrickRed, dotted] (A) -- (P);
    \draw[thick, BrickRed, dotted] (B) -- (P);
    \draw[thick, BrickRed, dotted] (C) -- (P);
    \draw[thick, RoyalPurple, dotted] (A) -- (Q);
    \draw[thick, RoyalPurple, dotted] (B) -- (Q);
    \draw[thick, RoyalPurple, dotted] (C) -- (Q);
    \draw[thick, dotted] (P) -- (D);
    \draw[thick, dotted] (P) -- (E);
    \draw[thick, dotted] (P) -- (F);
    \draw[thick] (A) -- (F);
    \draw[thick] (A) -- (E);
    \draw[thick] (F) -- (E);

    \node[above] at (A) {\footnotesize $A$};
    \node[left] at (B) {\footnotesize $B$};
    \node[right] at (C) {\footnotesize $C$};
    \draw[fill=BrickRed] (P) circle (0.05)
        node[left] {\footnotesize $P$};
    \draw[fill=RoyalPurple] (Q) circle (0.05)
        node[below] {\footnotesize $Q$};
    \draw[fill=BrickRed] (D) circle (0.05)
        node[right] {\footnotesize $D$};
    \draw[fill=BrickRed] (E) circle (0.05)
        node[right] {\footnotesize $E$};
    \draw[fill=BrickRed] (F) circle (0.05)
        node[left] {\footnotesize $F$};

    \pic[draw,"$\bullet$", angle radius=8mm, angle eccentricity=1.4]
        {angle=F--A--B};
    \pic[draw,"$\bullet$", angle radius=6mm, angle eccentricity=1.4]
        {angle=B--A--P};
    \pic[draw,"$\circ$", angle radius=8mm, angle eccentricity=1.4]
        {angle=P--A--Q};
    \pic[draw,"$\bullet$", angle radius=6mm, angle eccentricity=1.4]
        {angle=Q--A--C};
    \pic[draw,"$\bullet + \circ$", angle radius=8mm,
         angle eccentricity=1.6] {angle=C--A--E};
\end{tikzpicture}
\end{figure}

Because we do not yet know if there exists point $Q$,
let $l_A, l_B, l_C$ be the respective isogonal lines. Notice that
by construction, because $AF = AP = AE$, $l_A$ is the perpendicular
bisector of $FE$. Analogously, it is self-evident that $l_B$ and
$l_C$ are perpendicular bisector of $FD$ and $ED$ respectively.
In other words, the point of intersection of $l_A, l_B, l_C$
is the center of the circumcircle of $\triangle{DEF}$ and the center
of such circle and the points $D, E, F$ always exist. Thus,
$Q$ is well-defined.

I personally find this method more interesting because we can
still prove it without prior knowledge. Continuing with the proofs,
consider the following diagram for a special case.

\begin{figure}[H]
\centering
\begin{tikzpicture}[scale=0.8]
    \coordinate (A) at (2, 4);
    \coordinate (B) at (0, 0);
    \coordinate (C) at (6, 0);
    \coordinate (H) at (2, 2);
    \coordinate (O) at (3, 1);

    \draw[thick] (A) -- (B) -- (C) -- (A);
    \draw[thick, BrickRed, dotted] (A) -- (H);
    \draw[thick, BrickRed, dotted] (B) -- (H);
    \draw[thick, BrickRed, dotted] (C) -- (H);
    \draw[thick, RoyalPurple, dotted] (A) -- (O);
    \draw[thick, RoyalPurple, dotted] (B) -- (O);
    \draw[thick, RoyalPurple, dotted] (C) -- (O);
    \draw[thick, BrickRed] (A) -- (H);
    \draw[thick, BrickRed, dotted] (A) -- (2, 0);
    \draw[thick, RoyalPurple,] (A) -- (O);
    \draw[thick, RoyalPurple, dotted] (A) -- (4, 2);

    \fill[BrickRed, opacity=0.35]
        (A) -- (B) -- (2, 0) -- (A) -- cycle;
    \fill[RoyalPurple, opacity=0.35]
        (A) -- (O) -- (4, 2) -- (A) -- cycle;

    \node[above] at (A) {\footnotesize $A$};
    \node[left] at (B) {\footnotesize $B$};
    \node[right] at (C) {\footnotesize $C$};
    \draw[fill=BrickRed] (H) circle (0.05)
        node[left] {\footnotesize $H$};
    \draw[fill=RoyalPurple] (O) circle (0.05)
        node[below] {\footnotesize $O$};

    \pic[draw,"$\bullet$", angle radius=8mm, angle eccentricity=1.4]
        {angle=B--A--H};
    \pic[draw,"$\bullet$", angle radius=8mm, angle eccentricity=1.4]
        {angle=O--A--C};
    \pic[draw,"$\circ$", angle radius=8mm, angle eccentricity=1.4]
        {angle=H--B--A};
    \pic[draw,"$\circ$", angle radius=8mm, angle eccentricity=1.4]
        {angle=C--B--O};
    \pic[draw,"$\times$", angle radius=8mm, angle eccentricity=1.4]
        {angle=O--C--B};
    \pic[draw,"$\times$", angle radius=8mm, angle eccentricity=1.4]
        {angle=A--C--H};
\end{tikzpicture}
\end{figure}

Notice that because $\angle{AEF} = \angle{ABC}$, $\angle{BAH} =
\angle{OAC}$ by AA similarity. Similarly, the analogous argument
applies to the other angles.

\begin{important}
Points $H$ and $O$ are isogonal conjugates.
\end{important}

The concept of isogonal conjugate appears surprisingly a lot when
we discuss more advanced topics like harmonic range and spiral
similarity. This is it for this section! Let's discuss homothety
in the next section before getting into exciting topics like
harmonic range and spiral similarity.

\end{document}


